% ============================================================
%  ch14  第五讲 §5.4 融合积
%  原文行 2930–3074
% ============================================================

\subsection{融合积}\label{fusion}

Beilinson--Drinfeld Grassmann 簇的分解结构使我们能把卷积积 \eqref{Lus conv} 解释为融合积，由此可一并推出命题 \ref{fiber functor} 与命题 \ref{existence of comm}。

首先，用定义 $\on{P}_{L^+G}(\Gr_G)$ 的同一套手续（见 \S\ref{A:ind}），得到 $\Gr_{X^I}$ 上 $(L^+G)_{X^I}$-等变反常层的 Abel 范畴 $\on{P}_{(L^+G)_{X^I}}(\Gr_{X^I})$。仿照 Reich（见 \cite{Rei}），定义 $\Sat_{X^I}\subset \on{P}_{(L^+G)_{X^I}}(\Gr_{X^I})$ 为那些关于态射 $q_I:\Gr_{X^I}\to X^I$ 满足 ULA 条件的层构成的全子范畴（ULA 性质的综述见 \S\ref{ULA sheaf}）。由定理 \ref{III:factorization} 与定理 \ref{ULA perverse}，对每个 $\phi: J\to I$ 有函子
\begin{equation}\label{Sat restriction}
\Delta(\phi)^\bullet:=(\nu^\phi)^*[|I|-|J|]\cong(\nu^\phi)^![|J|-|I|]:\Sat_{X^J}\to\Sat_{X^I}.
\end{equation}

给定满射 $\phi:J\twoheadrightarrow I=\{1,2,\ldots,n\}$，有从卷积 Grassmann 簇到 Beilinson--Drinfeld Grassmann 簇的卷积映射
\[m_\phi:\Gr_{X^{J_1}}\tilde\times\cdots\tilde\times\Gr_{X^{J_n}}\to \Gr_{X^J}.\]
对 $\mA_i\in\Sat_{X^{J_i}}$，可用 \eqref{torsor E} 在 $\Gr_{X^{J_1}}\tilde\times\cdots\tilde\times\Gr_{X^{J_n}}$ 上构造「外扭曲积」$\mA_1\tilde\boxtimes\cdots\tilde\boxtimes\mA_n$。然后可定义「外卷积积」
\[\mA_1\boxstar \mA_2\boxstar \cdots\boxstar \mA_n:=(m_\phi)_!(\mA_1\tilde\boxtimes\cdots\tilde\boxtimes\mA_n).\]
关键陈述如下。

\begin{prop}\label{conv=fusion}~
\begin{enumerate}
\item 复形 $\mA_1\boxstar\mA_2\boxstar\cdots\boxstar\mA_n$ 属于 $\Sat_{X^{J}}$，因此有定义良好的外卷积积函子
\[\boxstar_{i\in I}:\Sat_{X^{J_1}}\times\cdots\times\Sat_{X^{J_n}}\to \Sat_{X^J},\quad (\mA_1,\ldots,\mA_n)\mapsto \mA_1\boxstar\cdots\boxstar\mA_n.\]

\item 有典范同构
\begin{equation}\label{eqfusion}
\mA_1\boxstar\cdots\boxstar\mA_n\cong j_{!*}((\mA_1\boxtimes\cdots\boxtimes\mA_n)|_{X^{(\phi)}}),
\end{equation}
其中 $j: X^{(\phi)}\to X^J$ 记开嵌入，而 $(\mA_1\boxtimes\cdots\boxtimes\mA_n)|_{X^{(\phi)}}$ 经同构 \eqref{BD factorization} 看作 $\Gr_{X^J}\times_{X^J}X^{(\phi)}$ 上的反常层。
\end{enumerate}
\end{prop}

\begin{proof}
首先注意若所有 $\mA_i$ 都 ULA，则 $\mA_1\tilde\boxtimes\cdots\tilde\boxtimes\mA_n$ 也 ULA。这是因为（在光滑拓扑下）局部上支集 $\on{Supp}(\mA_1\tilde\times\cdots\tilde\times\mA_n)$ 同构于乘积 $\on{Supp}(\mA_1)\times\cdots\times\on{Supp}(\mA_n)$。于是由定理 \ref{ULA prop} (3) 与 (4)' 即得。

再用定理 \ref{ULA prop} (2) 可知 $\mA_1\boxstar\cdots\boxstar\mA_n$ 关于 $\Gr_{X^J}\to X^J$ 是 ULA。此外它在 $X^{(\phi)}$ 上的限制反常，因为由引理 \ref{generic conv}，在同构 \eqref{BD factorization} 下
\[\mA_1\boxstar \mA_2\boxstar \cdots\boxstar \mA_n|_{X^{(\phi)}}=\mA_1\boxtimes\cdots\boxtimes\mA_n|_{X^{(\phi)}}.\]
故由定理 \ref{ULA perverse}，它反常，且 $\mA_1\boxstar\cdots\boxstar\mA_n\cong j_{!*}((\mA_1\boxtimes\cdots\boxtimes\mA_n)|_{X^{(\phi)}})$。最后，$(L^+G)_{X^J}$-等变性显然。
\end{proof}

下面推出若干推论（即下文 (i)--(iii)），特别是给出命题 \ref{fiber functor} 的另一证明与命题 \ref{existence of comm} 的证明。

(i) 结合 \eqref{Sat restriction}，可在每个 $\Sat_{X^J}$ 上定义内卷积积函子
\begin{equation}\label{internal convolution}
\ostar_{i\in I}: \times_{i\in I}\Sat_{X^J}\to \Sat_{X^J},\quad \mA_1\ostar\mA_2\ostar\cdots\ostar\mA_n:= \Delta^\bullet(\mA_1\boxstar\mA_2\boxstar\cdots\boxstar\mA_n),
\end{equation}
其中 $\Delta: X^I\to (X^I)^n$ 是对角嵌入，$\Delta^\bullet=\Delta^*[-(n-1)|I|]$。与 $\Sat_G$ 同样的理由表明每个 $\Sat_{X^J}$ 都是幺半范畴。

\begin{rmk}\label{triv group}
可以试着用同一公式 \eqref{internal convolution} 在整 $\on{P}_{(L^+G)_{X^J}}(\Gr_{X^J})$ 上定义内卷积积 $\ostar$。但下例说明 $\ostar$ 在这个更大的范畴上并不反常正合，这也是我们要求 ULA 条件的原因之一。

设 $G=\{e\}$ 平凡。则 $\on{P}_{L^+G}(\Gr_X)$ 就是 $X$ 上反常层的范畴，而 $\Sat_X$ 是 $X$ 上局部系（平移 $-1$ 度）构成的范畴。按定义
\[\mA_1\boxstar\mA_2=\mA_1\boxtimes\mA_2,\quad \mA_1\ostar\mA_2=\mA_1\otimes\mA_2[-1].\]
特别地，若 $\mA_1$ 或 $\mA_2$ 是摩天楼层，则 $\mA_1\ostar\mA_2$ 不反常。
\end{rmk}

注意有自然函子
\[\Sat_G\to \Sat_X,\quad \mA\mapsto \mA_X:=\Ql\tilde\boxtimes\mA[1],\]
其中 $\mA_X=\Ql\tilde\boxtimes\mA[1]$ 是 $\Gr_X=X\tilde\times\Gr$ 上的外扭曲积（见 \eqref{III: twisted constr}）。

\begin{lem}~
\begin{enumerate}
\item 函子 $\mA\mapsto \mA_X$ 完全忠实。
\item 对每个 $n$ 有典范同构
\[(\mA_1\star\cdots\star\mA_n)_X\cong (\mA_1)_X\ostar\cdots\ostar(\mA_n)_X.\]
特别地，函子 $\Sat_G\to \Sat_X,\ \mA\mapsto \mA_X$ 上有自然幺半结构。
\end{enumerate}
\end{lem}

\begin{proof}
注意图
\[X\tilde\times \on{Supp}(\mA)\leftarrow \hat X^+\times \on{Supp}(\mA)\to \on{Supp}(\mA)\]
中的两个映射都可写成纤维几何连通的平滑满射的逆向极限，故 (1) 由 \cite[Proposition 4.2.5]{BBD} 推出。(2) 由典范同构
\[X\tilde\times (\Gr\tilde\times\cdots\tilde\times\Gr)\cong \Gr_X\tilde\times\cdots\tilde\times\Gr_X|_\Delta\]
推出。
\end{proof}

\begin{rmk}\label{conv perverse}
本引理与命题 \ref{conv=fusion} 合起来给出命题 \ref{conv prod} 的证明。事实上它证明了更强的陈述（如 \cite{Ga} 中命题 6 之后的注所述），因为无需假设 $\mA_1$ 是 $L^+G$-等变的。
\end{rmk}

(ii) 注意 $(q_I)_*\mA$ 在 $X^I$ 上 ULA。令
\begin{equation}\label{relative coh}
\mH^*(\mA):=\bigoplus_i \mH^i(q_I)_*\mA[i].
\end{equation}
由定理 \ref{ULA prop} (1)，这是 $X^I$ 上的分次局部系。于是有典范同构
\begin{equation}\label{monoidal functor}
\mH^*(\mA_1\boxstar\cdots\boxstar\mA_n)\cong \mH^*(\mA_1)\boxtimes\cdots\boxtimes\mH^*(\mA_n),
\end{equation}
因为由 K\"unneth 公式它在 $X^{(\phi)}$ 上成立。于是得到到 $X^I$ 上分次局部系范畴的函子 $\mH^*: \Sat_{X^I}\to \on{Loc}^{\on{gr}}(X^I)$，它与外卷积积、张量积相容。因此 $\mH^*$ 关于内卷积积与张量积有自然幺半结构
\[\mH^*: (\Sat_{X^I},\ostar)\to (\on{Loc}^{\on{gr}}(X^I),\otimes).\]
对每个点 $i_x:\{x\}\to X$，取茎 $i_x^*$ 是从 $\on{Loc}^{\on{gr}}(X)$ 到分次向量空间范畴 $\on{Vect}^{\on{gr}}_{\Ql}$ 的对称幺半函子（两者都配以分次交换约束 $c_{\on{gr}}$）。虽然 $i_x^*$ 依赖点 $x$ 的选取，复合
\begin{equation}\label{decomp coh}
\Sat_G\to \Sat_X\stackrel{\mH}{\to} \on{Loc}^{\on{gr}}(X)\stackrel{i_x^*}{\to}\on{Vect}^{\on{gr}}_{\Ql}
\end{equation}
却不依赖，且由 proper 基变换典范同构于超上同调函子 $\on{H}^*$。于是也得到 $\on{H}^*$ 上的幺半结构。留给读者验证：它与命题 \ref{fiber functor} 中构造的幺半结构一致。

(iii) \eqref{eqfusion} 右端的表达式通常称为 $\mA_1,\cdots,\mA_n$ 的\textbf{融合积}。这一解释显出卷积积更多来自 $X^I$ 对称性的对称性。下面的简单陈述已足以证明命题 \ref{existence of comm}；更一般的讨论见注 \ref{symmetry on fusion}。

设 $\mA_1,\mA_2\in\Sat_X$。回忆有对合 $\Delta(\sigma):\Gr_{X^2}\cong\Gr_{X^2}$，见命题 \ref{III:Gr2}。由该命题与 \eqref{eqfusion} 可知存在典范同构
\[\Delta(\sigma)^*(\mA_1\boxstar\mA_2)\cong \mA_2\boxstar\mA_1\]
使下图交换
\[\begin{CD}
\sigma^*\mH^*(\mA_1\boxstar\mA_2)@>\cong>> \mH^*\Delta(\sigma)^*(\mA_1\boxstar\mA_2)@>\cong>>\mH^*(\mA_2\boxstar\mA_1)\\
@V\cong VV & & @VV\cong V\\
\sigma^*(\mH^*(\mA_1)\boxtimes\mH^*(\mA_2))@>c_{\on{gr}}>> \mH^*(\mA_2)\boxtimes\mH^*(\mA_1),
\end{CD}\]
其中 $c_{\on{gr}}$ 是（例如在 $(x,y)\in X^2$ 处）由
\[\begin{split}\sigma^*(\mH^*(\mA_1)\boxtimes\mH^*(\mA_2))_{(x,y)}=\mH^*(\mA_1)_y\otimes\mH^*(\mA_2)_x \\ \stackrel{c_{\on{gr}}}{\cong}\mH^*(\mA_2)_x\otimes \mH^*(\mA_1)_y=(\mH^*(\mA_2)\boxtimes\mH^*(\mA_1))_{(x,y)}\end{split}\]
给出的同构。

因此

\begin{cor}\label{existence of comm2}
存在典范同构 $c'_{\mA_1,\mA_2}:\mA_1\ostar\mA_2\cong\mA_2\ostar\mA_1$ 使下图交换
\[\begin{CD}
\mH^*(\mA_1\ostar\mA_2)@>c'_{\mA_1,\mA_2}>\cong>\mH^*(\mA_2\ostar\mA_1)\\
@V\cong VV@VV\cong V\\
\mH^*(\mA_1)\otimes\mH^*(\mA_2)@>c_{\on{gr}}>>\mH^*(\mA_2)\otimes\mH^*(\mA_1).
\end{CD}\]
\end{cor}

最后，因 $\Sat_G$ 是 $\Sat_X$ 的全子范畴，且超上同调函子 $\on{H}^*:\Sat_G\to \on{Vect}_{\Ql}$ 同构于 \eqref{decomp coh}，命题 \ref{existence of comm} 得证。

\begin{rmk}\label{symmetry on fusion}
更一般地，对每个 $\phi: J\twoheadrightarrow I$（$I$ 无次序）有函子
\[\underset{\phi}{\fstar}: \times_{i\in I}\Sat_{X^{J_i}}\to \Sat_{X^J},\quad ( \mA_{i}, i\in I)\mapsto \underset{\phi}{\fstar}\mA_i:=j_{!*}((\boxtimes_i \mF_i)|_{X^{(\phi)}}).\]
（右端关于 $q_I$ 是 ULA，这是命题 \ref{conv=fusion} 的推论。）于是由定理 \ref{III:factorization}，对每个 $K\stackrel{\phi}{\twoheadrightarrow} J\stackrel{\psi}{\twoheadrightarrow} I$ 有函子的典范同构
\begin{equation}\label{fusion prod}
\underset{\phi}{\fstar}= \underset{\psi\phi}{\fstar}(\times_{i\in I} \underset{\phi_i}{\fstar}): \times_{j} \Sat_{X^{K_j}}\to \Sat_{X^K},
\end{equation}
以及
\begin{equation}
\Delta(\phi)^\bullet \underset{\psi\phi}{\fstar}= \underset{\psi}{\fstar}(\times_{i\in I}\Delta(\phi_i)^\bullet): \times_{i}\Sat_{X^{K_i}}\to \Sat_{X^J}.
\end{equation}
由这些公式与 \cite[Proposition 1.5]{DM} 直接推出：$(\Sat_{X^I},\ostar)$ 有张量范畴结构，且 $\mH^*: (\Sat_{X^I},\ostar)\to (\on{Loc}^{\on{gr}}(X^I),\otimes)$ 有张量结构。
\end{rmk}

\begin{ex}
设 $G=\GL_n$。上述用融合积构造交换约束的方式与例 \ref{momega1BD} 合起来蕴含：在例 \ref{conv v.s. springer} 中把 $\Gr_{\omega_1,\ldots,\omega_1}\cap m^{-1}(t^\la)$ 等同于 Springer 纤维的该等同下，对称群 $S_m$ 在 $\Hom(V_\la, V_{\omega_1}^{\otimes m})$ 上的作用与 Springer 作用相差一个符号。特别地，由此得到 Schur--Weyl 对偶。注意该等同也是 Laumon 从 Springer 层构造自守层的出发点（参见 \cite{FGKV}）。
\end{ex}
